% ====================================================================
% Section: Background and Related Work (Revised)
% ====================================================================
\section{Background and Related Work}\label{sec:background}

This section establishes the theoretical groundwork and contextualizes the contributions of this thesis. Section \ref{subsec:rl_concepts} first introduces the foundational principles of Reinforcement Learning (RL) and the specific Deep Reinforcement Learning (DRL) algorithms central to this work. Section \ref{subsec:related_work_iot} then provides a review of related research in IoT security, situating this project within the broader academic landscape and highlighting its novel contributions.

% --- Subsection 2.1 ---
\subsection{Reinforcement Learning for Autonomous Systems}\label{subsec:rl_concepts}
Reinforcement Learning is a computational paradigm focused on goal-directed learning from interaction. An intelligent agent learns to make optimal decisions not by being shown labeled examples, but by exploring an environment and discovering which sequences of actions yield the greatest cumulative reward \cite{Sutton2018}. This learning process is formalized through the mathematical framework of Markov Decision Processes (MDPs).

\subsubsection{Markov Decision Processes}\label{subsubsec:mdp}
An MDP provides a formal model for sequential decision-making under uncertainty. The core principle is the \textbf{Markov Property}, which asserts that the current state provides all necessary information to make an optimal decision, rendering the history of prior states irrelevant \cite{Sutton2018}. An MDP is defined as a five-tuple $(S, A, P, R, \gamma)$:
\begin{itemize}
    \item \textbf{S}: A finite set of states, representing all possible configurations of the environment. In this thesis, a state $s \in S$ is a complex observation composed of current network traffic features, a history of recent states and actions, and predictions from the LSTM attack model.
    \item \textbf{A}: A finite set of actions. For this work, $A$ consists of four discrete defensive actions: \{`monitor`, `rate-limit`, `block IP`, `shutdown service`\}.
    \item \textbf{P}: The state transition probability function, $P(s' | s, a) = P(S_{t+1}=s' | S_t=s, A_t=a)$, defining the dynamics of the environment.
    \item \textbf{R}: The reward function, $R(s, a, s')$, which provides the immediate numerical feedback for transitioning from $s$ to $s'$ via action $a$.
    \item \textbf{$\gamma$}: The discount factor ($0 \leq \gamma \leq 1$), which prioritizes immediate rewards over future ones.
\end{itemize}
The agent's objective is to learn a \textbf{policy}, $\pi(a|s)$, which maps states to a probability distribution over actions, to maximize the expected discounted return.

\subsubsection{Value Functions and Bellman Equations}\label{subsubsec:bellman}
To determine the quality of states and actions, RL uses value functions. The \textbf{state-value function}, $V^\pi(s)$, is the expected return from state $s$ while following policy $\pi$:
\begin{equation}\label{eq:state_value_function}
V^\pi(s) = \mathbb{E}_\pi \left[ \sum_{k=0}^{\infty} \gamma^k R_{t+k+1} \Big| S_t = s \right]
\end{equation}
The \textbf{action-value function} (or Q-function), $Q^\pi(s, a)$, is the expected return after taking action $a$ from state $s$ and then following policy $\pi$:
\begin{equation}\label{eq:action_value_function}
Q^\pi(s,a) = \mathbb{E}_\pi \left[ \sum_{k=0}^{\infty} \gamma^k R_{t+k+1} \Big| S_t = s, A_t = a \right]
\end{equation}
These functions adhere to the recursive relationships defined by the Bellman equations. The \textbf{Bellman optimality equation} for the Q-function is central to many RL algorithms and is expressed as:
\begin{equation}\label{eq:bellman_optimality}
Q^*(s, a) = \mathbb{E}_{s' \sim P(\cdot|s,a)} \left[ R(s,a,s') + \gamma \max_{a' \in A} Q^*(s', a') \right]
\end{equation}
Solving this equation yields the optimal Q-function, $Q^*(s, a)$, from which the optimal policy can be derived by selecting the action with the highest value in each state.

\subsubsection{Deep Reinforcement Learning Algorithms}\label{subsubsec:drl_algorithms}
For environments with large or continuous state spaces, like an IoT network, tabular methods are intractable. Deep Reinforcement Learning (DRL) overcomes this by using deep neural networks to approximate the value functions or policies \cite{Chen2021}. This thesis implements and benchmarks three seminal DRL algorithms that represent the primary families of RL.

\paragraph{Deep Q-Networks (DQN)}
DQN is a value-based algorithm that uses a neural network to approximate the action-value function, $Q(s, a; \theta) \approx Q^*(s, a)$ \cite{Mnih2015}. To stabilize learning, DQN introduces two critical mechanisms:
\begin{itemize}
    \item \textbf{Experience Replay:} Experiences $(s_t, a_t, r_t, s_{t+1})$ are stored in a large replay buffer. The network is trained on mini-batches randomly sampled from this buffer, which breaks temporal correlations in the data.
    \item \textbf{Fixed Target Network:} A separate, periodically updated target network is used to generate stable target values for the Bellman equation, preventing oscillations during training. The loss function is typically the Mean Squared Bellman Error (MSBE):
    \begin{equation}\label{eq:dqn_loss}
    L(\theta) = \mathbb{E}_{(s,a,r,s') \sim D} \left[ \left( \underbrace{(r + \gamma \max_{a'} Q(s', a'; \theta^-))}_{\text{Target Q-value}} - \underbrace{Q(s, a; \theta)}_{\text{Predicted Q-value}} \right)^2 \right]
    \end{equation}
    where $\theta^-$ are the parameters of the fixed target network.
\end{itemize}

\paragraph{Proximal Policy Optimization (PPO)}
PPO is a state-of-the-art actor-critic algorithm that directly optimizes a parameterized policy (the actor) while using a value function (the critic) to reduce variance \cite{Schulman2017}. Its key innovation is a clipped surrogate objective function that constrains the size of policy updates, ensuring greater training stability. The objective function is:
\begin{equation}\label{eq:ppo_clip}
L^{CLIP}(\theta) = \hat{\mathbb{E}}_t \left[ \min \left( r_t(\theta)\hat{A}_t, \text{clip}(r_t(\theta), 1-\epsilon, 1+\epsilon)\hat{A}_t \right) \right]
\end{equation}
where $r_t(\theta)$ is the probability ratio between the new and old policies, $\hat{A}_t$ is the estimated advantage, and $\epsilon$ is a hyperparameter that constrains the policy change. This approach prevents destructive policy updates and is known for its reliability and performance.

\paragraph{Advantage Actor-Critic (A2C)}
A2C is a synchronous variant of the foundational Asynchronous Advantage Actor-Critic (A3C) algorithm \cite{Mnih2016}. Like PPO, it is an actor-critic method, but it uses a simpler, unconstrained policy update mechanism. It typically employs multiple parallel workers to collect experience simultaneously. These experiences are then aggregated to compute a single, stable gradient update for both the actor (policy) and the critic (value function). While often faster to converge than PPO on simpler problems, it can be less stable.

All three algorithms are implemented in this thesis using Stable Baselines3's `MultiInputPolicy`, which is designed to handle the `Dict` observation space provided by our custom IoT environment.

% --- Subsection 2.2 ---
\subsection{Related Work in AI-Driven IoT Security}\label{subsec:related_work_iot}
The application of artificial intelligence to IoT security is an active and rapidly evolving field of research. Early work focused on applying classical machine learning models to intrusion detection, which has since paved the way for more sophisticated, adaptive systems based on Deep Reinforcement Learning.

\subsubsection{Machine Learning Approaches for Intrusion Detection}
Early approaches to intelligent IoT security leveraged supervised and unsupervised machine learning for building Intrusion Detection Systems (IDS). Supervised models, such as Support Vector Machines (SVM) and Random Forests, were trained on labeled datasets to classify network flows as either benign or malicious \cite{AlGaradi2020}. While effective at identifying known attack patterns, these models struggle with novel or \textbf{zero-day attacks} not seen during training, a limitation they share with traditional signature-based IDS \cite{Khraisat2019}. Furthermore, they typically operate as passive classifiers, identifying threats without the capacity for automated, sequential defensive actions.

\subsubsection{The Shift to Active Defense with DRL}
The limitations of static ML models motivated a shift towards DRL to create autonomous agents capable of active, closed-loop cyber-defense. Several notable works have explored this domain, providing a foundation upon which this thesis builds.

Some research has focused on using RL for deception and intelligence gathering. The work of Guan et al. (2023) in \textbf{HoneyIoT}, for example, presents an adaptive high-interaction honeypot that uses RL to learn an attacker's behavior and dynamically alter its services to keep the attacker engaged and reveal their methods \cite{Guan2023}. While this is a powerful application of RL, its primary purpose is reconnaissance, whereas our system is designed as a direct defense mechanism for a live network.

Other works have targeted specific layers of the IoT ecosystem. \textbf{IoTWarden}, proposed by Alam et al. (2024), uses a DQN agent to mitigate attacks in trigger-action IoT platforms (like IFTTT) by identifying and disabling malicious rules or "applets" in real-time \cite{Alam2024}. This approach is highly effective for its targeted domain but is confined to the application layer. Similarly, the work by Nguyen et al. (2021) proposes a Federated DRL system for traffic monitoring in Software-Defined Networks (SDN) for IoT, focusing on efficient data collection and analysis rather than active defense \cite{Nguyen2021}. Our research addresses a more fundamental challenge by operating at the \textbf{network layer}, analyzing features derived from raw traffic to defend against a broad range of network-level attacks such as DDoS, DoS, and Spoofing, which are prevalent in the CICIoT2023 dataset.

\subsubsection{Positioning the Thesis: Bridging the Research Gap}
While previous works have successfully applied DRL to specific aspects of IoT security, a gap remains for a comprehensive, network-level defense system that is simultaneously:
\begin{itemize}
    \item \textbf{Proactive and Adaptive:} Uniquely combining a data-driven predictive model (an LSTM) with a DRL-based agent to create a closed loop from threat anticipation to autonomous action.
    \item \textbf{Grounded in Realistic, Modern Data:} Validated on the large-scale and contemporary \textbf{CICIoT2023 dataset}, ensuring that the learned policies are relevant to modern, real-world attack vectors.
    \item \textbf{Network-Layer Focused:} Designed to defend against a broad range of fundamental network attacks by analyzing traffic-flow features.
    \item \textbf{Rigorously Benchmarked:} Providing a thorough, empirical comparison of modern value-based (DQN) and actor-critic (PPO, A2C) algorithms for this specific proactive defense task.
\end{itemize}
This thesis directly addresses this gap by designing, implementing, and evaluating a complete framework that integrates these critical components, thereby advancing the development of resilient, intelligent systems capable of defending the vast and ever-expanding IoT landscape.